\chapter{Continuidad y equivalencia topológica}\label{chap:05}
\section{Controlar la imagen de puntos próximos}
Una función continua transforma aproximaciones válidas en el dominio en aproximaciones válidas en el codominio. Hay tres lenguajes para expresarlo: bolas y tolerancias, imágenes inversas de abiertos y sucesiones. La lectura 5 de \cite{valencia} conecta estos lenguajes. Es esencial distinguir la \emph{imagen} $f(A)$ de la \emph{imagen inversa} $f^{-1}(B)$: la primera se aplica a subconjuntos del dominio, la segunda a subconjuntos del codominio, y la notación de imagen inversa no exige que $f$ sea invertible.

\begin{definicion}\label{ch05:continua}\index{continuidad!en un punto}
Sean $(X,d)$ y $(Y,\rho)$ espacios métricos. Una función $f:X\to Y$ es continua en $x\in X$ si para toda vecindad $V$ de $f(x)$, el conjunto $f^{-1}(V)$ es vecindad de $x$.
\end{definicion}

\begin{teorema}[Criterio de tolerancias]\label{ch05:epsdelta}
La función $f$ es continua en $x$ si y solo si
\[
\forall\eps>0\ \exists\delta>0\ \forall u\in X:
\quad d(x,u)<\delta\ \Longrightarrow\ \rho(f(x),f(u))<\eps.
\]
\end{teorema}
\begin{proof}
Si $f$ es continua en $x$, la imagen inversa de la vecindad $B_Y(f(x),\eps)$ contiene alguna bola $B_X(x,\delta)$. La inclusión de bolas es exactamente la implicación del enunciado. Recíprocamente, sea $V$ una vecindad de $f(x)$ y elijamos $\eps>0$ con $B_Y(f(x),\eps)\subset V$. La condición proporciona $\delta>0$ con $B_X(x,\delta)\subset f^{-1}(B_Y(f(x),\eps))\subset f^{-1}(V)$, así que esta última es vecindad de $x$.
\end{proof}
El orden de los cuantificadores explica el significado: primero se prescribe una precisión en las imágenes; después se elige una tolerancia en el dominio que sirve para \emph{todos} los puntos suficientemente próximos. La tolerancia puede depender del punto $x$; exigir una sola tolerancia para todos los puntos conduce más adelante a continuidad uniforme.

\begin{teorema}[Criterios globales]\label{ch05:criterios}\index{continuidad!criterios topologicos@criterios topológicos}
Para $f:X\to Y$, son equivalentes:
\begin{enumerate}[label=(\roman*)]
\item $f^{-1}(U)$ es abierto en $X$ para todo abierto $U$ en $Y$.
\item Para todo $A\subset X$, $f(\cl A)\subset\cl{f(A)}$.
\item $f^{-1}(F)$ es cerrado en $X$ para todo cerrado $F$ en $Y$.
\item $f$ es continua en cada punto de $X$.
\end{enumerate}
Cuando se cumplen, se dice que $f$ es continua en $X$.
\end{teorema}
\begin{proof}
\emph{(i)$\Rightarrow$(ii).} Si $x\in\cl A$ y $U$ es abierto con $f(x)\in U$, entonces $f^{-1}(U)$ es un abierto que contiene a $x$. Por ello intersecta a $A$, y la imagen de un punto de esa intersección pertenece a $U\cap f(A)$. Como toda bola alrededor de $f(x)$ es uno de esos abiertos, $f(x)\in\cl{f(A)}$.

\emph{(ii)$\Rightarrow$(iii).} Sea $F$ cerrado y $E=f^{-1}(F)$. La inclusión $f(E)\subset F$ implica $\cl{f(E)}\subset F$. Si $x\in\cl E$, por (ii), $f(x)\in f(\cl E)\subset\cl{f(E)}\subset F$, de modo que $x\in E$. Así $E$ es cerrado.

\emph{(iii)$\Rightarrow$(i).} Para $U$ abierto, $Y\setminus U$ es cerrado y
\[
X\setminus f^{-1}(U)=f^{-1}(Y\setminus U)
\]
es cerrado. Por tanto $f^{-1}(U)$ es abierto.

\emph{(i)$\Rightarrow$(iv).} Si $V$ es vecindad de $f(x)$, contiene una bola abierta $U$ centrada en $f(x)$. Por (i), $f^{-1}(U)$ es abierto y contiene a $x$; contiene entonces una bola alrededor de $x$. Como $f^{-1}(U)\subset f^{-1}(V)$, esta última es vecindad de $x$.

\emph{(iv)$\Rightarrow$(i).} Sea $x\in f^{-1}(U)$, con $U$ abierto. Este $U$ es vecindad de $f(x)$, así que $f^{-1}(U)$ es vecindad de $x$. Lo es para cada uno de sus puntos y, por tanto, es abierto.
\end{proof}
No se afirma en (ii) igualdad. Por ejemplo, $f(t)=\arctan t$ es continua en $\R$, pero para $A=\R$ su imagen es $(-\pi/2,\pi/2)$ y la clausura de esa imagen contiene además los dos extremos.

\section{Ejemplos y reglas de construcción}
Toda función cuyo dominio lleva la topología discreta es continua, porque cualquier imagen inversa es un subconjunto abierto del dominio. La identidad de un espacio en sí mismo es continua, ya que sus imágenes inversas no cambian los conjuntos.

Si $f:X\to Y$ y $g:Y\to Z$ son continuas, también lo es $g\circ f$: para un abierto $U\subset Z$,
\[
(g\circ f)^{-1}(U)=f^{-1}(g^{-1}(U)),
\]
y primero la continuidad de $g$ y luego la de $f$ garantizan apertura. Se está usando la preimagen de conjuntos, no una inversa funcional de $f$ o $g$.

Una función $L$-Lipschitz, es decir, una función con $\rho(f(x),f(y))\le Ld(x,y)$, es continua. Si $L>0$, basta tomar $\delta=\eps/L$; si $L=0$, todas las imágenes coinciden y cualquier $\delta>0$ sirve. Así, la función $x\mapsto d(x,A)$ para $A\ne\varnothing$ es continua por el ejercicio 7 del capítulo~\ref{chap:03}, y la norma $v\mapsto\norm v$ es continua porque
\[
|\norm u-\norm v|\le\norm{u-v}.
\]
Esta última desigualdad se prueba aplicando la triangular a $u=(u-v)+v$ y luego intercambiando $u,v$.

La inclusión $Y\hookrightarrow X$ de un subespacio es continua: la preimagen de un abierto ambiental es su intersección con $Y$, abierta relativamente. Una restricción $f|_Y$ de una función continua es continua por composición con esa inclusión. Además, si una función toma valores en un subespacio $Z\subset Y$, es continua como función a $Z$ si y solo si lo es como función a $Y$. En una dirección se compone con la inclusión; en la otra, para $U=Z\cap G$ abierto en $Z$, se tiene $f^{-1}(U)=f^{-1}(G)$ porque toda imagen pertenece a $Z$. Estos hechos incorporan el resultado sobre subespacios que aparece repetido en la lectura original.

\begin{teorema}[Criterio secuencial]\label{ch05:secuencial}\index{continuidad!criterio secuencial}
Una función entre espacios métricos es continua en $x$ si y solo si $x_n\to x$ implica $f(x_n)\to f(x)$.
\end{teorema}
\begin{proof}
Supongamos continuidad y fijemos $\eps>0$. Elegimos $\delta>0$ según el teorema~\ref{ch05:epsdelta}. Como $x_n\to x$, hay $N$ tal que $d(x_n,x)<\delta$ para $n\ge N$. Entonces $\rho(f(x_n),f(x))<\eps$ para esos índices.

Para el recíproco usamos la negación exacta de continuidad: existe $\eps_0>0$ tal que para \emph{todo} $\delta>0$ existe $u\in X$ con $d(u,x)<\delta$ y $\rho(f(u),f(x))\ge\eps_0$. Tomando $\delta=1/n$, elegimos $x_n$ con esas propiedades. Entonces $x_n\to x$, pero ninguna imagen entra en la bola $B(f(x),\eps_0)$. Esto contradice el criterio secuencial. La desigualdad de la negación es $\ge\eps_0$, y no es necesario reemplazarla por una estricta.
\end{proof}

\section{Funciones abiertas, cerradas y homeomorfismos}
\begin{definicion}\index{funcion@función!abierta}\index{funcion@función!cerrada}\index{homeomorfismo}
Una función $f:X\to Y$ es abierta si $f(U)$ es abierto para cada abierto $U\subset X$, y cerrada si $f(F)$ es cerrado para cada cerrado $F\subset X$. Es un homeomorfismo si es biyectiva, continua y su inversa $f^{-1}:Y\to X$ es continua. En este caso se escribe $X\cong Y$.
\end{definicion}
Estas condiciones sobre imágenes directas son diferentes de continuidad. La identidad de $\R$ con su métrica discreta hacia $\R$ usual es continua, pero no es abierta, pues la imagen de $\{0\}$ no es abierta, ni cerrada, pues la imagen del conjunto discreto cerrado $\Q$ no es cerrada en la recta usual. La identidad en la dirección contraria es abierta y cerrada, porque todo subconjunto del codominio discreto es abierto y cerrado; sin embargo, no es continua.

\begin{ejemplo}\label{ch05:homeorecta}
La función $h:\R\to(-1,1)$ dada por $h(t)=t/(1+|t|)$ es un homeomorfismo. Su valor tiene módulo menor que $1$. Para $s\in(-1,1)$, definamos $k(s)=s/(1-|s|)$. Como $|k(s)|=|s|/(1-|s|)$,
\[
h(k(s))=\frac{s/(1-|s|)}{1+|s|/(1-|s|)}=s.
\]
Del mismo modo, $1-|h(t)|=1/(1+|t|)$ implica $k(h(t))=t$. Por tanto son funciones inversas y ambas son continuas por las reglas elementales de cálculo: los denominadores son positivos en sus respectivos dominios. Esto demuestra biyectividad y continuidad de las dos direcciones, no solamente una semejanza geométrica.
\end{ejemplo}

\section{Ejercicios resueltos}
\begin{ejercicio}[Lectura 5, ejercicio 1]
Demostrar que cada una de las siguientes condiciones equivale a continuidad de $f:X\to Y$:
\begin{enumerate}[label=(\roman*)]
\item Para todo $A\subset Y$, $f^{-1}(\operatorname{int}A)\subset\operatorname{int}(f^{-1}(A))$.
\item Para todo $A\subset Y$, $\cl{f^{-1}(A)}\subset f^{-1}(\cl A)$.
\end{enumerate}
En (ii), se corrige la letra $B$, no definida en el original, por $A$.
\end{ejercicio}
\begin{solucion}
\begin{enumerate}[label=(\roman*)]
\item Si $f$ es continua, $f^{-1}(\operatorname{int}A)$ es abierto y está contenido en $f^{-1}(A)$. El interior de un conjunto contiene a todos sus subconjuntos abiertos; de ahí la inclusión. Si la condición vale y $U$ es abierto, entonces $\operatorname{int}U=U$ y $f^{-1}(U)\subset\operatorname{int}(f^{-1}(U))$. La inclusión inversa es automática, de modo que $f^{-1}(U)$ es abierto y $f$ continua.
\item Si $f$ es continua, $f^{-1}(\cl A)$ es cerrado y contiene a $f^{-1}(A)$; por ello contiene a su clausura. Recíprocamente, para un cerrado $F$ tenemos $\cl F=F$, y la condición da $\cl{f^{-1}(F)}\subset f^{-1}(F)$. Por tanto todas las preimágenes de cerrados son cerradas y $f$ es continua.
\end{enumerate}
\end{solucion}

\begin{ejercicio}[Lectura 5, ejercicio 2]
Demostrar que una función es continua en el espacio si y solo si es continua en cada punto.
\end{ejercicio}
\begin{solucion}
Las dos direcciones están desarrolladas en el teorema~\ref{ch05:criterios}. El mecanismo es local: la preimagen de un abierto $U$ es abierta exactamente cuando contiene una bola alrededor de cada $x$ con $f(x)\in U$. La continuidad en ese $x$ produce esa bola, y la continuidad global garantiza que la preimagen de una bola alrededor de $f(x)$ es un abierto que contiene a $x$. Así se recupera en cada dirección la definición requerida.
\end{solucion}

\begin{ejercicio}[Lectura 5, ejercicio 3]
Probar que $f$ es continua si y solo si las preimágenes de todas las bolas abiertas son abiertas. Dar una función no continua cuyas preimágenes de todas las bolas cerradas sean cerradas.
\end{ejercicio}
\begin{solucion}
La necesidad sigue de que las bolas abiertas son abiertas. Para la suficiencia, todo abierto $U$ del codominio es una unión $\bigcup_iB_i$ de bolas abiertas. Como las preimágenes preservan uniones,
\[
f^{-1}(U)=\bigcup_i f^{-1}(B_i),
\]
que es abierto por hipótesis.

Para la segunda parte, tomemos $X=\R$ con la métrica usual, $Y=\R$ con la discreta y $f(t)=t$. Una bola cerrada de radio $r<1$ en $Y$ es un singleton, y una de radio $r\ge1$ es todo $Y$. Sus preimágenes son singletons o $\R$, todos cerrados en $X$. Pero $\{0\}$ es abierto en $Y$ y su preimagen no es abierta en $X$; así $f$ no es continua. El argumento con bolas abiertas no se transfiere a bolas cerradas, porque una unión arbitraria de cerrados no tiene por qué ser cerrada.
\end{solucion}

\begin{ejercicio}[Lectura 5, ejercicio 4]
Sean $f,g:X\to\R$ continuas. Probar que (i) $\{x:f(x)<g(x)\}$ es abierto y (ii) $\{x:f(x)\le g(x)\}$ es cerrado. Se corrige el dominio de $g$ a $X$, necesario para que ambas expresiones estén definidas.
\end{ejercicio}
\begin{solucion}
\begin{enumerate}[label=(\roman*)]
\item Si $f(x)<g(x)$, sea $\eta=(g(x)-f(x))/3>0$. La continuidad de ambas funciones permite elegir $\delta>0$ tal que $d(x,y)<\delta$ implique $|f(y)-f(x)|<\eta$ y $|g(y)-g(x)|<\eta$. Entonces
\[
f(y)<f(x)+\eta<g(x)-\eta<g(y).
\]
La bola está contenida en el conjunto, que es abierto.
\item El complemento es $\{x:g(x)<f(x)\}$, abierto por (i) aplicado con los papeles intercambiados. Por tanto el conjunto original es cerrado.
\end{enumerate}
\end{solucion}

\begin{ejercicio}[Lectura 5, ejercicio 5]\index{extension@extensión!por cero}\index{frontera}
Sea $A\subset X$ cerrado y $f:A\to\R$ continua, con $f=0$ en la frontera $\partial A=\cl A\cap\cl{X\setminus A}$. Demostrar que la extensión $g:X\to\R$, igual a $f$ en $A$ y a cero fuera de $A$, es continua.
\end{ejercicio}
\begin{solucion}
Distinguimos tres posiciones del punto $x$.

Si $x\notin A$, como $A$ es cerrado existe una bola alrededor de $x$ contenida en $X\setminus A$. Allí $g$ es constante y es continua en $x$.

Si $x\in\operatorname{int}A$, existe $r>0$ con $B_X(x,r)\subset A$. Dado $\eps>0$, la continuidad relativa de $f$ da $\delta_0>0$ tal que $y\in A$ y $d(x,y)<\delta_0$ implican $|f(y)-f(x)|<\eps$. Tomando $\delta=\min\{r,\delta_0\}$ se obtiene la misma desigualdad para $g$ y todo $y\in X$ con $d(x,y)<\delta$.

Finalmente, si $x\in A\setminus\operatorname{int}A$, todo entorno de $x$ intersecta al complemento de $A$, así que $x\in\partial A$ y $g(x)=f(x)=0$. Por continuidad relativa existe $\delta>0$ con $|f(y)|<\eps$ para $y\in A\cap B(x,\delta)$. Para $y$ de esa bola fuera de $A$, $g(y)=0$; para los puntos dentro de $A$, $|g(y)|=|f(y)|<\eps$. La misma tolerancia sirve para ambos lados de la frontera. Por tanto $g$ es continua en todos los puntos.
\end{solucion}

\begin{ejercicio}[Lectura 5, ejercicio 6: corrección del recíproco]\index{soporte}
Sea $S=\supp(f)=\cl{\{x:f(x)\ne0\}}$. Examinar la afirmación «$f$ es continua si y solo si $f=0$ en $\partial S$» y establecer una versión válida.
\end{ejercicio}
\begin{solucion}
La implicación directa es verdadera. Si $x\in\partial S$ y $f(x)\ne0$, por continuidad hay una bola en la que $|f(y)-f(x)|<|f(x)|/2$, así que $|f(y)|>|f(x)|/2>0$. Esa bola está contenida en $S$, lo que convierte a $x$ en punto interior de $S$ y contradice $x\in\partial S$.

El recíproco es falso. En $\R$, tomemos $f(x)=1$ para $x<0$ y $f(x)=2$ para $x\ge0$. La función no es continua en $0$, pero nunca se anula, de modo que $S=\R$ y $\partial S=\varnothing$. La condición en la frontera se satisface de forma vacía y no controla el salto interior.

Una versión exacta es: $f$ es continua en $X$ si y solo si $f|_S$ es continua en el subespacio $S$ y $f=0$ en $\partial S$. La necesidad se sigue de la continuidad de las restricciones y de la prueba anterior. Para la suficiencia, $S$ es cerrado por definición, $f$ es cero en $X\setminus S$ y el ejercicio 5 aplica a la función continua $f|_S$. Su extensión por cero es precisamente $f$.
\end{solucion}

\begin{ejercicio}[Lectura 5, ejercicio 7]
Probar: (i) $f$ es abierta si y solo si las imágenes de las bolas abiertas son abiertas; (ii) si $f$ es biyectiva, es abierta si y solo si es cerrada; (iii) construir una función suprayectiva abierta que no sea cerrada y una cerrada que no sea abierta.
\end{ejercicio}
\begin{solucion}
\begin{enumerate}[label=(\roman*)]
\item La necesidad es inmediata porque una bola abierta es un abierto. Para la suficiencia, escribimos un abierto $U$ como unión de bolas $B_i$; entonces $f(U)=\bigcup_if(B_i)$ es una unión de abiertos.
\item Para una biyección se cumple $f(X\setminus A)=Y\setminus f(A)$. Si $f$ es abierta y $F$ cerrado, el conjunto $Y\setminus f(F)=f(X\setminus F)$ es abierto, así que $f(F)$ es cerrado. Si $f$ es cerrada y $U$ abierto, el conjunto $Y\setminus f(U)=f(X\setminus U)$ es cerrado, de modo que $f(U)$ es abierto. La identidad de complementos utiliza tanto inyectividad como suprayectividad.
\item La proyección $\pi:\R^2\to\R$, $\pi(x,y)=x$, es suprayectiva y abierta. Si $a\in\pi(U)$ con $U$ abierto, existe $b$ con $(a,b)\in U$ y un disco de radio $r>0$ centrado en $(a,b)$ contenido en $U$. Entonces $(a-r,a+r)\subset\pi(U)$, pues $(t,b)$ pertenece a ese disco cuando $|t-a|<r$. Sin embargo, el conjunto $F=\{(x,y):xy=1\}$ es cerrado, como preimagen de $\{1\}$ por la función continua producto, y $\pi(F)=\R\setminus\{0\}$ no es cerrado.

Para una función incluso suprayectiva y cerrada que no sea abierta, tomemos $h:\R\to[0,\infty)$, $h(t)=\max\{t,0\}$. Si $F$ es cerrado, su imagen es $F\cap[0,\infty)$, con el punto $0$ añadido cuando $F$ intersecta $(-\infty,0)$. En ambos casos es un cerrado de $[0,\infty)$. Pero $h((-2,-1))=\{0\}$ no es abierto en $[0,\infty)$: toda bola relativa alrededor de $0$ contiene puntos positivos.
\end{enumerate}
\end{solucion}

\begin{ejercicio}[Lectura 5, ejercicio 8]\index{pegado!sobre abiertos}
Sea $(A_i)_{i\in I}$ una cubierta abierta de $X$ y sean $f_i:A_i\to Y$ continuas que coinciden en todas las intersecciones. Probar que la función obtenida pegándolas es continua.
\end{ejercicio}
\begin{solucion}
Para $x\in X$, elijamos $i$ con $x\in A_i$ y definamos $g(x)=f_i(x)$. Si también $x\in A_j$, la hipótesis da $f_i(x)=f_j(x)$; por tanto la definición no depende de la elección. Cualquier función con las restricciones prescritas tiene necesariamente este valor en cada $x$, de modo que es única.

Sea $U$ abierto en $Y$. Entonces
\[
g^{-1}(U)=\bigcup_{i\in I}f_i^{-1}(U).
\]
Cada preimagen es abierta en $A_i$ por continuidad. Como $A_i$ es abierto en $X$, el ejercicio 2 del capítulo~\ref{chap:04} garantiza que también es abierta en $X$. Su unión es abierta y el criterio global prueba continuidad de $g$. No se requiere que la cubierta sea finita, porque las uniones arbitrarias de abiertos preservan apertura.
\end{solucion}

\begin{ejercicio}[Lectura 5, ejercicio 9]
Demostrar que ser homeomorfos es una relación de equivalencia entre espacios métricos.
\end{ejercicio}
\begin{solucion}
La identidad de $X$ es una biyección continua cuya inversa es ella misma; esto prueba reflexividad. Si $h:X\to Y$ es un homeomorfismo, su inversa $h^{-1}:Y\to X$ es biyectiva y continua y tiene inversa continua $h$; esto prueba simetría. Si $h:X\to Y$ y $k:Y\to Z$ son homeomorfismos, $k\circ h$ es una biyección continua y su inversa es $h^{-1}\circ k^{-1}$, continua por composición. Así se prueba transitividad. Cuando se habla de la clase de todos los espacios, «relación de equivalencia» se entiende en ese sentido de clases; sobre cualquier colección que sea un conjunto se obtiene la relación de equivalencia usual.
\end{solucion}

\section{Elegir el criterio adecuado}
El criterio de tolerancias permite obtener cotas concretas; el de abiertos simplifica composiciones y pegados; el secuencial permite detectar discontinuidades mediante aproximaciones explícitas. Un homeomorfismo conserva esta estructura de continuidad en ambas direcciones. No tiene por qué conservar distancias, acotación ni propiedades que dependan de la métrica concreta.
